\documentclass{article}
\usepackage{enumitem}
\usepackage{amsmath}
\begin{document}

\title{Chapter 5: Morphology of Flowering Plants}
\date{}
\maketitle

\begin{enumerate}[label=Q\arabic*.]

\item Which of the following is a modified stem and NOT a root?
\\(A) Carrot
\\(B) Sweet potato
\\(C) Ginger (rhizome)
\\(D) Turnip

\item The tap root system is characteristic of:
\\(A) Monocots
\\(B) Dicots
\\(C) Both monocots and dicots equally
\\(D) Pteridophytes

\item Which type of phyllotaxy is seen in \textit{Calotropis}?
\\(A) Alternate
\\(B) Opposite
\\(C) Whorled
\\(D) Decussate

\item A leaf with reticulate venation and two cotyledons in its seed belongs to:
\\(A) Monocots
\\(B) Dicots
\\(C) Gymnosperms
\\(D) Pteridophytes

\item Which of the following is a correct description of an epigeal germination?
\\(A) Cotyledons remain below soil, plumule emerges
\\(B) Cotyledons are pushed above the soil by elongation of hypocotyl
\\(C) Both cotyledons and plumule emerge simultaneously
\\(D) Radicle emerges before cotyledons

\item Assertion: The onion bulb is a modified underground stem.\\
Reason: It has nodes, internodes, and axillary buds covered by fleshy leaf bases.
\\(A) Both true, Reason is correct explanation
\\(B) Both true, Reason is not correct explanation
\\(C) Assertion true, Reason false
\\(D) Assertion false, Reason true

\item Racemose inflorescence differs from cymose inflorescence in that:
\\(A) The main axis terminates in a flower in racemose
\\(B) The main axis has unlimited growth in racemose
\\(C) Older flowers are at the top in racemose
\\(D) Lateral branches are absent in racemose

\item The thalamus contributes to fruit formation in:
\\(A) True fruits like mango
\\(B) False fruits like apple
\\(C) Both true and false fruits
\\(D) Only aggregate fruits

\item A superior ovary with axile placentation is found in:
\\(A) \textit{Dianthus}
\\(B) \textit{Lemon}
\\(C) \textit{Poppy}
\\(D) \textit{Marigold}

\item Which of the following is a aggregate fruit?
\\(A) Mango
\\(B) Strawberry
\\(C) Pineapple
\\(D) Banana

\item The formula of a flower written as $\oplus \; \wp \; K_{(5)} \; C_{(5)} \; A_5 \; G_{(2)}$ indicates:
\\(A) Actinomorphic, bisexual, 5 sepals fused, 5 petals fused, 5 stamens free, 2 carpels fused, superior ovary
\\(B) Zygomorphic, unisexual, 5 sepals free, 5 petals fused, 5 stamens fused, 2 carpels free
\\(C) Actinomorphic, bisexual, 5 sepals free, 5 petals fused, 5 stamens free, 2 carpels fused, inferior ovary
\\(D) Zygomorphic, bisexual, 5 sepals fused, 5 petals free, 5 stamens free, 2 carpels fused

\item Which type of root modification is seen in \textit{Cuscuta} (dodder)?
\\(A) Pneumatophores
\\(B) Haustorial roots (parasitic roots)
\\(C) Prop roots
\\(D) Stilt roots

\item The calyx, corolla, androecium, and gynoecium are collectively called:
\\(A) Perianth
\\(B) Whorls of a flower
\\(C) Accessory and essential organs
\\(D) Floral appendages

\item Match floral features with their families:\\
1. Epipetalous stamens $\rightarrow$ (a) Solanaceae\\
2. Diadelphous stamens $\rightarrow$ (b) Fabaceae\\
3. Zygomorphic corolla $\rightarrow$ (c) Fabaceae\\
4. Valvate aestivation $\rightarrow$ (d) Malvaceae
\\(A) 1-a, 2-b, 3-c, 4-d
\\(B) 1-b, 2-a, 3-d, 4-c
\\(C) 1-c, 2-d, 3-a, 4-b
\\(D) 1-d, 2-c, 3-b, 4-a

\item Vivipary, the germination of seeds while still attached to the parent plant, is an adaptation seen in:
\\(A) Desert plants
\\(B) Mangrove plants in tidal marshes
\\(C) Epiphytes
\\(D) Alpine plants

\end{enumerate}
\end{document}
